\documentclass[10pt]{article}
\usepackage[margin=0.85in]{geometry}
\usepackage{amsmath,amssymb}
\usepackage[colorlinks=true,linkcolor=blue,urlcolor=blue]{hyperref}
\usepackage{xcolor}
\setlength{\parskip}{3pt}\setlength{\parindent}{0pt}

\newcommand{\Z}{\mathbb{Z}}
\newcommand{\F}{\mathbb{F}}
\newcommand{\Q}{\mathbb{Q}}
\DeclareMathOperator{\ord}{ord}
\DeclareMathOperator{\lcm}{lcm}

\begin{document}

\begin{center}
{\LARGE\bfseries Research memo: the arithmetic-geometric direction}\\[3pt]
{\small PRIMES-Conjectures project --- AG track, branch
\texttt{arith-geom-interpretation} --- October 21--22, 2026\\
Full development: \texttt{doc/arith-geom.tex} (v1.0, 16 pp);
standalone Popovych note: \texttt{doc/popovych-note.tex} (5 pp);
results log: \texttt{doc/arith-geom-findings.md}}
\end{center}

\textbf{TL;DR.} Reading Agrawal's conjecture through algebraic groups turned
the project's empirical dead ends (A11, N1, A10) into theorems, produced the
first quantified evidence on Popovych's conjecture beyond search ranges,
closed the elliptic-curve escape route, and gave the negative-result paper
its structural spine: a single \emph{stabilizer-index theorem} explaining
why every variant of the attack bottoms out at the same twin-smooth wall.
Everything is unconditional except one named Artin-type hypothesis ---
\emph{generic sharpness} (the support bounds are attained; $0$ exceptions
across $196$ sampled large primes, and a theorem in the function-field
mirror); the restriction to Frobenius-matching mechanisms is a separate,
stated scope, not a hypothesis.

\section*{1. The dictionary}

The AKS congruence $(X-1)^n \equiv X^n - 1 \pmod{n, X^r-1}$ is a
\emph{Frobenius-lift test}: $\Z[X]/(X^r{-}1)$ is the $\Lambda$-ring of
$\mu_r$ with Adams operations $\psi^k(X) = X^k$, and $n$ is prime iff
$\psi^n$ is a Frobenius lift at $n$ (Wilkerson); AKS checks the lift
condition on the single section $X - 1$. The exception clause
$n^2 \equiv 1 \pmod r$ says exactly that the fake Frobenius $\sigma_n$ lies
in $\mathrm{Gal}(\Q(\zeta_r)/\Q(\zeta_r)^+) = \{1, \text{conj}\}$ --- the
\emph{stabilizer} of the section's divisor. This ``clause $=$ stabilizer''
principle recurs in every result below.

On the counterexample side, the Lenstra--Pomerance conditions
$p - 1 \mid n - 1$, $p + 1 \mid n + 1$ say that $n$ \emph{impersonates the
Frobenius of $p$ on both one-dimensional algebraic tori over $\F_p$} (split:
order $p-1$; nonsplit norm-one: order $p+1$). Twin-smoothness of
$(\frac{p-1}{2}, \frac{p+1}{2})$ is torus-order smoothness --- the same
condition behind the B-SIDH/SQIsign parameter hunts, which is why the
project's reliance on the Costello--Naehrig corpora was never a coincidence.

\section*{2. The four results}

\textbf{(AG2) The torus dimension-count theorem --- A11 is now a
corollary.} For $d = \ord_r(p)$, the unique involution of $(\Z/r)^\times$
being $-1$ forces the exact \emph{inversion collapse}
$(X-1)^{p^{d/2}-1} = -X^{-1}$ (even $d$), so the control order
$\rho_r(p)$ lives on cyclotomic levels $e \mid d'$, $d' = d/2$ or $d$, and
the smoothness burden is the torus dimension
$\sum_{e \mid d'} \varphi(e) = d'$. Validity ($\ord_r(n) > 2$) forces
$d \ge 3$, so \textbf{burden $\ge 2$ for every AKS modulus $r$, with the
minimum $2$ attained exactly at $d = 4$, levels $\{1,2\}$: the double
(twin-smooth) condition} --- realizable iff $4 \mid r - 1$, and V\'a\v{n}a's
$\rho_5(p) \mid 10(p^2{-}1)$ is the special case. The burden formula
reproduces the A11 empirical table at all 11 moduli (value and minimizing
$d$); verified on 67 sample primes, drawn and implemented independently of
the A11 scan. ``Pick a friendlier $r$'' is closed
structurally, not just empirically.

\textbf{(AG5) Popovych is strictly harder --- first quantified evidence.}
The collapse above is a property of $X - 1$, whose zero $\{1\}$ is
Galois-stable. The second Popovych section $X + 2$ has \emph{trivial}
stabilizer --- $(\zeta^\gamma{+}2)/(\zeta{+}2)$ is not even a unit of
$\Z[\zeta_r]$ for $\gamma \ne 1$ (disjoint ideal supports: norms $r$ vs
$(2^r{+}1)/3$) --- so no collapse exists at any depth and the pair pays full
burden $\min\{d : d \mid r{-}1, d \ge 3\} \ge 3$, equal to $4$ at Popovych's
own $r = 5$: pool primes would need $p-1$, $p+1$ \emph{and} $p^2+1$
simultaneously smooth. At the frozen harvest spec (pool primes
${\approx}51$ bits) that extra ${\approx}102$-bit level is a Dickman tax
$\rho(102/\log_2 B)$: $3.9\times10^{-11}$ over the solver base $B = 1259$,
still $1.8\times10^{-7}$ over the generous $B' = 12590$ ---
emptying the measured pool. Bonus structure: the second congruence's
$X{=}1$ component is exactly the base-3 Fermat test $3^n \equiv 3$, and
mod $3$ the two sections coincide.

\textbf{(AG5b, Oct 22 addendum) Does the ``2'' matter? No --- but ``not
$1$'' is subtly incomplete.} For any integer $c \notin \{0, \pm 1\}$ the
trivial-stabilizer argument runs verbatim for $X + c$ and the full burden
jump holds, with $c$ entering only through the norm
$\Phi_r(-c) = (c^r{+}1)/(c{+}1)$ (finitely many excluded $p$, with
$p \mid c(c{+}1)$), the excluded modulus $r \mid c{+}1$, and the level-one
component --- the base-$(c{+}1)$ Fermat test. What must be avoided is
exactly the rational torsion $\{\pm 1\}$ as the section's zero $-c$:
$c = -1$ is the first section itself, and the innocent-looking $c = +1$
fails \emph{outright}. Its zero $-1$ is inversion-fixed, so the collapse
survives ($\ord(X{+}1) \mid r(p^{d/2}{-}1)$ at even $d$, measured
$22/22$); worse, $\zeta{+}1 = (\zeta^2{-}1)/(\zeta{-}1)$ with both
numerator and denominator values of $X - 1$ at primitive roots, whose
orders divide $\rho_r(p)$ --- so the mechanism congruences plus the
Carmichael condition imply $T(+1, n, r)$ \emph{automatically} (proved:
Prop.~8 of the note / Prop.~5.9 of the full document; measured: $48/48$
synthetic mechanism instances pass $T(+1)$, $0/48$ pass $T(+2)$). The
hypothetical ``$\{X{-}1, X{+}1\}$ conjecture'' --- Popovych's statement
with $2$ replaced by $1$ --- is therefore \emph{mechanism-equivalent} to
Agrawal's: formally stronger-looking, but every counterexample the known
engine could ever build defeats both simultaneously. In stabilizer terms
$U(X \pm 1) = \{\pm 1\}$ while $U(X + c) = \{1\}$ otherwise: $c = 2$ is
simply the smallest constant that genuinely blocks the engine ($c = -2$
works equally, with Mersenne norm $2^r - 1$).

\textbf{(AG3) The impersonation hierarchy --- our targets have a name, and
the wall has a shape.} Define cells by (level set, Frobenius powers); at
level 2 the power's parity is a sign. Then: Carmichael $= (\{1\}, 0)$,
infinite (AGP); Lucas--Carmichael $= (\{2\}, 1)$, infinite (Wright 2018);
the L--P targets $= (\{1,2\}, (0,1))$ are the \emph{Williams numbers}
(Williams 1977, Echi) --- open even for nonemptiness, none below $2^{64}$
(McIntosh--Mitra 2015). Yet Howe's \emph{rigid order-2 Carmichael number}
$443{,}372{,}888{,}629{,}441 = 17{\cdot}31{\cdot}41{\cdot}43{\cdot}89
{\cdot}97{\cdot}167{\cdot}331$ occupies the \emph{same level set} with
all-even parities and exists at $4.4 \times 10^{14}$ (re-verified).
\textbf{Sign splitting, not depth, is the wall}: rigid cells are one-sided
AGP problems (single target $1$); the Williams parities $(0,1)$ force
targets $(+1, -1)$ that conflict on shared prime powers, forcing the
coprime side-split whose measured collapse is N1/A10. Validity pins Agrawal
to the sign-split cell; sharp datum: Howe's $n \equiv 1 \pmod 5$ does not
even pass $T(-1,n,5)$ --- rigidity misses the $\mu_{2r}$ torsion the
mod-80 conditions exist to fix.

\textbf{(AG4) The elliptic axis is closed --- the stabilizer-index
theorem.} A canonical elliptic analogue of $X \mapsto X^n$ exists only for
CM curves ($\Lambda$-structure iff CM, Borger--Gurney). There Frobenius
acts on an $r$-eigenline by a scalar $\lambda$, curve automorphisms give
exact collapses (e.g.\ $x([i]S) = -x(S)$ on $y^2 = x^3 - x$), and
everything unifies:
\[
\textbf{burden} \;=\; [\langle\lambda\rangle :
\langle\lambda\rangle \cap U] \;\ge\; 2 \quad
\text{for valid } \lambda \notin U,
\]
$U$ the unit-stabilizer of the section ($\{\pm1\}$ for $X{-}1$ $\to$ AG2;
$\{1\}$ for $X{+}2$ $\to$ AG5; $\mu_4/\mu_6$ for CM $x$-coordinates). The
tantalizing elliptic burden-\emph{1} cells ($\lambda \equiv \pm i$, section
order collapsing to $p - 1$ alone) are exactly the exception clause --- the
elliptic mirror of rigid Carmichaels --- and valid burden-2 cells demand
twin-smoothness again, with sign-split targets. Measured on
$y^2 = x^3 - x$: valid-cell floors $3$ ($r{=}13$) and $2$ ($r{=}17$),
matching the index formula; elliptic $\ge$ torus burden at all 11 moduli.
Gate AG-G2 \textbf{not passed}; the negative result extends to the
CM-elliptic axis.

\section*{3. Epistemic status}

Proved outright: the collapse identities, the support bounds, the
independence/stabilizer lemmas, the index formula, the burden
minimizations. One ingredient is heuristic (Artin-type \emph{generic
sharpness}: the support bounds are attained), empirically supported with
$0$ exceptions across $196$ sampled large primes --- and it is a
\emph{theorem} in the function-field mirror via Weil's RH for curves
(AG6 appendix), locating the residual difficulty over $\Z$ in the
unconditional-equidistribution gap. All results are mechanism-scoped
(Frobenius-matching constructions --- every published counterexample
approach); four deterministic verification scripts
(\texttt{code/ag\{2,3,4,5\}\_*.py}) with committed JSON back every claim.

\section*{4. Where this leaves the program}

\begin{itemize}\itemsep1pt
\item \textbf{Negative-result paper}: written and revised to v1.1 (branch
\texttt{negative-result-writeup}, 12 pp): \S5 now \emph{leads} with the
stabilizer-index theorem (with a per-section table: section / $U$ /
minimum burden / realizability) and derives the modulus, Popovych and
elliptic results as corollaries; the one conditional ingredient is named
in the abstract; \S7 adds two constructive reframings --- the index as a
\emph{design metric} (security parameter) for AKS-style tests, and the
unconditional function-field floor --- per external review feedback.
The obstruction is structural, not parametric: no modulus, no second
section, no CM curve escapes the floor, whose minimum is exactly the
twin-smooth sign-split cell that N1/A10 measured dead.
\item \textbf{Standalone Popovych note}: now written
(\texttt{doc/popovych-note.tex}, 5 pp) --- self-contained, conditionality
in the abstract, the non-unit lemma, the burden theorem with the $r = 5$
computation and comparison table, the Dickman estimate with explicit
parameters, and the AG5b addendum (Prop.~8: the $X{+}1$ variant is
mechanism-equivalent to Agrawal's). First quantified evidence beyond
search ranges; ready to circulate.
\item \textbf{The only live constructive direction} remains option (b): a
fundamentally new solver for $0/1$ subset-products in ${\sim}24$k-bit
smooth-exponent groups at identity-density ${\sim}0.004$. The AG results
sharpen its stakes: a breakthrough would unblock the torus route only ---
Popovych and elliptic variants stay strictly harder.
\item \textbf{Open problems surfaced}: unconditional sharpness
(equidistribution); existence of Williams numbers at all (independent
interest; the sign-split lens suggests why the AGP machine hasn't reached
them); whether any side-separating polynomial family reaches an admissible
local cell (old Q6.2, still open).
\end{itemize}

\end{document}
